% ==================== Art-Book Editorial Preamble for AI Series Volume 1 ====================
\usepackage[top=24mm,bottom=24mm,left=23mm,right=23mm,headheight=22pt,headsep=7mm]{geometry}
\usepackage{amsmath,amssymb,mathtools,amsthm}
\usepackage{booktabs,tabularx,array,multirow,longtable}
\usepackage{graphicx}
\usepackage{xcolor}
\usepackage{tikz}
\usepackage{pgfplots}
\pgfplotsset{compat=1.18}
\usetikzlibrary{arrows.meta,positioning,shapes.geometric,calc,decorations.pathreplacing,backgrounds}
\usepackage[most]{tcolorbox}
\usepackage{fancyhdr}
\usepackage{listings}
\usepackage{enumitem}
\usepackage{hyperref}

% ==================== Curated Gallery Color Palette ====================
\definecolor{ink}{HTML}{1B2430}        % Deep charcoal-slate ink
\definecolor{teal}{HTML}{0B5555}       % Refined verdigris teal
\definecolor{tealDark}{HTML}{083B3B}   % Deep gallery spine teal
\definecolor{tealLight}{HTML}{F2F8F8}  % Delicate ice-teal parchment
\definecolor{tealLine}{HTML}{A3C6C5}   % Soft architectural rule
\definecolor{tealMist}{HTML}{DCEBEB}   % Watermark chapter numeral tint
\definecolor{navy}{HTML}{1D3B5E}       % Muted Prussian indigo
\definecolor{navyLight}{HTML}{F2F5FA}  % Soft indigo tint
\definecolor{ochre}{HTML}{9E4A22}      % Burnished terracotta-ochre
\definecolor{gold}{HTML}{B97A24}       % Fine burnished amber-gold accent
\definecolor{ochreLight}{HTML}{FCF5F0} % Warm parchment tint
\definecolor{plum}{HTML}{5B3A6B}       % Muted editorial plum
\definecolor{plumLight}{HTML}{F8F4FB}  % Soft plum tint
\definecolor{slate}{HTML}{526071}      % Muted architectural gray
\definecolor{codebg}{HTML}{F7F9FB}     % Clean light code canvas
\definecolor{codeborder}{HTML}{D4DCE5}

\hypersetup{
  unicode=true,
  colorlinks=true,
  urlcolor=teal,
  linkcolor=teal,
  citecolor=teal,
  pdftitle={天书中的智能 · 第1卷：Python/NumPy编程、AI线性代数与多元微积分},
  pdfauthor={天书中的智能：人工智能的七卷手记}
}

% ==================== Editorial Chapter Openers & Typography ====================
\makeatletter
\newcommand{\chapWatermarkNum}{%
  \@ifundefined{inappendix}{\ifnum\value{chapter}<10 0\fi\thechapter}{\thechapter}%
}
\makeatother

\ctexset{
  contentsname={目\quad 录},
  chapter={
    name={第,章},
    number=\arabic{chapter},
    format=\LARGE\sffamily\bfseries\color{tealDark},
    nameformat=\small\sffamily\bfseries\color{gold},
    aftername=\hfill{\fontsize{38}{38}\selectfont\sffamily\bfseries\color{tealMist}\chapWatermarkNum}\par\vspace{-4mm},
    aftertitle=\par\vspace{2.8mm}{\noindent\color{gold}\rule{26mm}{1.5pt}}{\color{tealLine}\rule{\dimexpr\linewidth-26mm\relax}{0.55pt}}\par,
    beforeskip=6pt,
    afterskip=10pt
  },
  section={
    format=\Large\sffamily\bfseries\color{teal},
    beforeskip=14pt plus 3pt minus 2pt,
    afterskip=6pt plus 1pt
  },
  subsection={
    format=\large\sffamily\bfseries\color{ink},
    beforeskip=10pt plus 2pt minus 1pt,
    afterskip=4pt plus 1pt
  },
  subsubsection={
    format=\normalsize\sffamily\bfseries\color{navy},
    beforeskip=7pt,
    afterskip=2pt
  }
}
\renewcommand{\lstlistingname}{代码}
\renewcommand{\figurename}{图}
\renewcommand{\tablename}{表}

% Artisanal Chapter Epigraph macro (quote + attribution)
\newcommand{\chapEpigraph}[2]{%
  \vspace{-2mm}%
  \begin{tcolorbox}[
    enhanced,
    colback=ochreLight!45!white,
    colframe=ochreLight!45!white,
    borderline west={1.8pt}{0pt}{gold},
    boxrule=0pt,
    arc=0.8mm,
    left=4mm,right=3.5mm,top=2mm,bottom=2mm,
    before skip=0pt,after skip=8pt
  ]
    {\kaishu\normalsize\color{ink!85}“#1”}\par\vspace{1mm}%
    \hfill{\footnotesize\sffamily\color{slate}——\,#2}%
  \end{tcolorbox}%
}

% Exercise tier header macro
\newcommand{\exTier}[2]{%
  \par\vspace{3.8mm}\noindent
  \colorbox{tealLight}{\strut\;\sffamily\bfseries\footnotesize\color{tealDark}#1\;}\;%
  {\small\sffamily\bfseries\color{slate}#2}\par\nopagebreak\vspace{1.2mm}%
}
\newcommand{\exHint}[1]{\par\smallskip\noindent{\footnotesize\sffamily\color{navy}\textbf{【提示】}\,#1}}

\setlength{\parindent}{0pt}
\setlength{\parskip}{5pt plus 1.5pt minus 0.5pt}
\linespread{1.26}
\setlist[itemize]{leftmargin=1.8em,topsep=3pt,itemsep=2.5pt,parsep=1pt}
\setlist[enumerate]{leftmargin=2.0em,topsep=3pt,itemsep=2.5pt,parsep=1pt}
\setlength{\tabcolsep}{5.5pt}
\renewcommand{\arraystretch}{1.22}

% ==================== Running Headers & Footers ====================
\pagestyle{fancy}
\fancyhf{}
\fancyhead[L]{\small\sffamily\color{slate}天书中的智能 \;\textcolor{gold}{$\cdot$}\; 第 1 卷}
\fancyhead[R]{\small\sffamily\color{slate}\nouppercase{\leftmark}}
\fancyfoot[C]{\small\sffamily\color{slate}\textcolor{gold}{$\diamond$}\quad\textcolor{tealDark}{\thepage}\quad\textcolor{gold}{$\diamond$}}
\renewcommand{\headrulewidth}{0.4pt}
\renewcommand{\headrule}{\hbox to\headwidth{\color{tealLine}\leaders\hrule height \headrulewidth\hfill}}
\renewcommand{\chaptermark}[1]{\markboth{第 \thechapter\ 章\quad #1}{}}

\fancypagestyle{plain}{
  \fancyhf{}
  \fancyfoot[C]{\small\sffamily\color{slate}\textcolor{gold}{$\diamond$}\quad\textcolor{tealDark}{\thepage}\quad\textcolor{gold}{$\diamond$}}
  \renewcommand{\headrulewidth}{0pt}
}

% ==================== Theorem & Pedagogical Box Environments ====================
\theoremstyle{definition}
\newtheorem{theorem}{定理}[chapter]
\newtheorem{lemma}[theorem]{引理}
\newtheorem{proposition}[theorem]{命题}
\newtheorem{corollary}[theorem]{推论}
\newtheorem{definition}[theorem]{定义}
\newtheorem{example}[theorem]{例}
\newtheorem{remark}[theorem]{注记}
\renewcommand{\proofname}{\textbf{\sffamily\color{teal}证明}}

\newtcolorbox{thmBox}[1][]{
  enhanced,breakable,
  colback=tealLight,
  colframe=tealLine,
  borderline west={2.4pt}{0pt}{teal},
  boxrule=0.45pt,
  arc=1.5mm,
  left=4mm,right=3.8mm,top=2.8mm,bottom=2.8mm,
  before skip=6.5pt,after skip=6.5pt,
  #1
}

\newtcolorbox{defBox}[1][]{
  enhanced,breakable,
  colback=navyLight,
  colframe=navy!30,
  borderline west={2.4pt}{0pt}{navy},
  boxrule=0.45pt,
  arc=1.5mm,
  left=4mm,right=3.8mm,top=2.8mm,bottom=2.8mm,
  before skip=6.5pt,after skip=6.5pt,
  #1
}

\newtcolorbox{insightBox}[1]{
  enhanced,breakable,
  colback=white,
  colframe=navy!35,
  borderline west={2.4pt}{0pt}{gold},
  coltitle=navy,
  fonttitle=\sffamily\bfseries\small,
  attach title to upper={\quad},
  title={【直觉与下游坐标】#1},
  boxrule=0.45pt,
  arc=1.2mm,
  left=4mm,right=3.8mm,top=2.6mm,bottom=2.6mm,
  before skip=6.5pt,after skip=6.5pt
}

\newtcolorbox{recallBox}[1]{
  enhanced,breakable,
  colback=codebg,
  colframe=codeborder,
  borderline west={2.2pt}{0pt}{slate},
  coltitle=slate,
  fonttitle=\sffamily\bfseries\small,
  attach title to upper={\quad},
  title={【前置速忆】#1},
  boxrule=0.4pt,
  arc=1.2mm,
  left=3.8mm,right=3.5mm,top=2.4mm,bottom=2.4mm,
  before skip=5.5pt,after skip=5.5pt
}

\newtcolorbox{failBox}[1]{
  enhanced,breakable,
  colback=ochreLight,
  colframe=ochre!35,
  borderline west={2.4pt}{0pt}{ochre},
  coltitle=ochre,
  fonttitle=\sffamily\bfseries\small,
  attach title to upper={\quad},
  title={【反例、数值陷阱与排错】#1},
  boxrule=0.45pt,
  arc=1.2mm,
  left=4mm,right=3.8mm,top=2.6mm,bottom=2.6mm,
  before skip=6.5pt,after skip=6.5pt
}

\newtcolorbox{extBox}[1]{
  enhanced,breakable,
  colback=plumLight,
  colframe=plum!30,
  borderline west={2.4pt}{0pt}{plum},
  coltitle=plum,
  fonttitle=\sffamily\bfseries\small,
  attach title to upper={\quad},
  title={【选读拓展 · 初读可跳过】#1},
  boxrule=0.45pt,
  arc=1.2mm,
  left=4mm,right=3.8mm,top=2.6mm,bottom=2.6mm,
  before skip=6.5pt,after skip=6.5pt
}

\newtcolorbox{chapIntro}[1]{
  enhanced,
  colback=tealLight,
  colframe=tealLine,
  borderline west={2.8pt}{0pt}{tealDark},
  boxrule=0.45pt,
  arc=1.5mm,
  left=4.2mm,right=4mm,top=3mm,bottom=3mm,
  before skip=2pt,after skip=10pt,
  title={\sffamily\bfseries\small\color{tealDark}#1},
  colbacktitle=tealLight,
  attach boxed title to top left={yshift=-1mm,xshift=2mm},
  boxed title style={empty,size=minimal}
}

% ==================== Code Listing Style ====================
\lstdefinestyle{pycode}{
  language=Python,
  basicstyle=\footnotesize\ttfamily\color{ink},
  keywordstyle=\bfseries\color{teal},
  commentstyle=\itshape\color{slate},
  stringstyle=\color{ochre},
  backgroundcolor=\color{codebg},
  frame=single,
  rulecolor=\color{codeborder},
  framesep=5pt,
  xleftmargin=4pt,
  xrightmargin=4pt,
  numbers=left,
  numberstyle=\tiny\ttfamily\color{slate!80},
  numbersep=6pt,
  columns=fullflexible,
  keepspaces=true,
  showstringspaces=false,
  breaklines=true,
  aboveskip=6.5pt,
  belowskip=6.5pt
}
\lstset{style=pycode}

% ==================== Math Shortcuts ====================
\newcommand{\R}{\mathbb{R}}
\newcommand{\C}{\mathbb{C}}
\newcommand{\N}{\mathbb{N}}
\newcommand{\E}{\mathbb{E}}
\newcommand{\col}{\operatorname{col}}
\newcommand{\row}{\operatorname{row}}
\newcommand{\nul}{\operatorname{null}}
\newcommand{\rank}{\operatorname{rank}}
\newcommand{\diag}{\operatorname{diag}}
\newcommand{\tr}{\operatorname{tr}}
\newcommand{\spanop}{\operatorname{span}}
\newcommand{\argmin}{\operatorname*{arg\,min}}
\newcommand{\argmax}{\operatorname*{arg\,max}}
\newcommand{\cond}{\kappa}
\newcommand{\norm}[1]{\left\|#1\right\|}
\newcommand{\abs}[1]{\left|#1\right|}
\newcommand{\inner}[2]{\left\langle #1,\, #2 \right\rangle}
\newcommand{\had}{\odot}

\makeatletter
\setlength{\@fpsep}{16pt}
\makeatother
